% Configuração tipográfica inspirada nas normas ABNT para trabalhos acadêmicos.
% A classe article e o padrão tipográfico seguem a base existente da equipe.
% O conteúdo é organizado especificamente para o avulso de princípios.

\usepackage{fontspec}
\defaultfontfeatures{Ligatures=TeX}
\IfFileExists{fontes/TimesNewRoman.ttf}{%
  \setmainfont[Path=fontes/,UprightFont=TimesNewRoman.ttf,BoldFont=TimesNewRoman-Bold.ttf,ItalicFont=TimesNewRoman-Italic.ttf,BoldItalicFont=TimesNewRoman-BoldItalic.ttf]{TimesNewRoman}%
}{\setmainfont{texgyretermes-regular.otf}[BoldFont=texgyretermes-bold.otf,ItalicFont=texgyretermes-italic.otf,BoldItalicFont=texgyretermes-bolditalic.otf]}
\IfFileExists{fontes/Arial.ttf}{%
  \setsansfont[Path=fontes/,UprightFont=Arial.ttf,BoldFont=Arial-Bold.ttf,ItalicFont=Arial-Italic.ttf,BoldItalicFont=Arial-BoldItalic.ttf]{Arial}%
}{\setsansfont{texgyreheros-regular.otf}[BoldFont=texgyreheros-bold.otf,ItalicFont=texgyreheros-italic.otf,BoldItalicFont=texgyreheros-bolditalic.otf]}
\IfFileExists{fontes/lmmono10-regular.otf}{\setmonofont[Path=fontes/]{lmmono10-regular.otf}}{\setmonofont{lmmono10-regular.otf}}

\usepackage[brazilian]{babel}
\usepackage{geometry}
\geometry{top=3cm,bottom=2cm,left=3cm,right=2cm}
\usepackage{microtype}
\usepackage{setspace}
\onehalfspacing

\usepackage{graphicx}
\usepackage{pdflscape}
\usepackage[table]{xcolor}
\usepackage{tikz}
\usetikzlibrary{arrows.meta,positioning,calc}
\usepackage{booktabs}
\usepackage{tabularx}
\usepackage{longtable}
\usepackage{array}
\usepackage{ragged2e}
\usepackage{caption}
\usepackage{float}
\usepackage{fancyhdr}
\usepackage{etoolbox}
\usepackage[authoryear,round]{natbib}
\usepackage{hyperref}

\definecolor{cinzaescuro}{HTML}{222222}
\definecolor{cinzamedio}{HTML}{666666}
\definecolor{cinzaclaro}{HTML}{F1F1F1}
\definecolor{linhacinza}{HTML}{B7B7B7}

\hypersetup{
  colorlinks=false,
  pdfborder={0 0 0},
  pdftitle={Princípios de solução para a proteção de louças sanitárias},
  pdfauthor={Equipe 2}
}

\setlength{\parindent}{1.25cm}
\setlength{\parskip}{0pt}
\setlength{\emergencystretch}{2em}
\clubpenalty=10000
\widowpenalty=10000
\displaywidowpenalty=10000

\makeatletter
\renewcommand\section{\@startsection{section}{1}{0pt}{1.8\baselineskip}{0.7\baselineskip}{\normalfont\normalsize\bfseries}}
\renewcommand\subsection{\@startsection{subsection}{2}{0pt}{1.4\baselineskip}{0.45\baselineskip}{\normalfont\normalsize\bfseries}}
\renewcommand\subsubsection{\@startsection{subsubsection}{3}{0pt}{1.2\baselineskip}{0.35\baselineskip}{\normalfont\normalsize\bfseries}}
\setlength{\@fptop}{0pt}
\makeatother

\captionsetup{
  font=small,
  labelfont=bf,
  justification=justified,
  singlelinecheck=false,
  labelsep=endash
}
\captionsetup[table]{position=above,skip=6pt}

\renewcommand{\arraystretch}{1.22}
\newcolumntype{Y}{>{\RaggedRight\arraybackslash}X}
\newcolumntype{L}[1]{>{\RaggedRight\arraybackslash}p{#1}}
\newcommand{\fonte}[1]{\par\vspace{2pt}\noindent{\footnotesize Fonte: #1}}
\newcommand{\notaquadro}[1]{\par\vspace{2pt}\noindent{\footnotesize Nota: #1}}
\newcommand{\pendente}[1]{\par\noindent{\small\textbf{Ponto a validar:} \textit{#1}}\par}
\newcommand{\etapafutura}[1]{\par\noindent{\small\textit{#1}}\par}

% Importa linhas sem inserir comandos de arquivo dentro do alinhamento da tabela.
\makeatletter
\newcommand{\linhasquadro}[1]{\@@input #1 }
\makeatother

\AtBeginEnvironment{itemize}{\setlength{\itemsep}{0.2em}\setlength{\topsep}{0.4em}}
\AtBeginEnvironment{enumerate}{\setlength{\itemsep}{0.2em}\setlength{\topsep}{0.4em}}

\pagestyle{fancy}
\fancyhf{}
\fancyhead[R]{\small\thepage}
\renewcommand{\headrulewidth}{0pt}
\renewcommand{\footrulewidth}{0pt}
\setlength{\headheight}{15pt}

\newenvironment{resumorelatorio}
  {\begin{singlespace}\small\noindent}
  {\par\end{singlespace}}

% Extensões específicas do avulso de princípios.
\usepackage{xurl}
\urlstyle{same}
\setlength{\bibsep}{0.5\baselineskip}
\newcommand{\campocaso}[2]{\par\noindent\textbf{#1} #2\par}
\newcommand{\celulaopcao}[4]{%
  \begin{minipage}[t]{\linewidth}\vspace{0pt}%
  \RaggedRight\textbf{#1} #2\par\vspace{1mm}%
  {\centering\includegraphics[width=0.91\linewidth,height=14.5mm,keepaspectratio]{#3}\par}%
  \vspace{0.6mm}{\fontsize{8}{9}\selectfont #4\par}%
  \end{minipage}}
\newcommand{\celulafuncao}[5]{%
  \begin{minipage}[t]{\linewidth}\vspace{0pt}%
  \RaggedRight\textbf{#1}\par #2\par\vspace{2mm}%
  {\fontsize{8.5}{10}\selectfont C1: #3\par C2: #4\par C3: #5\par}%
  \end{minipage}}
